\documentclass[14 pt]{extarticle}

	\usepackage[frenchb]{babel}
	\usepackage[utf8]{inputenc}  
	\usepackage[T1]{fontenc}
	\usepackage{amssymb}
	\usepackage[mathscr]{euscript}
	\usepackage{stmaryrd}
	\usepackage{amsmath}
	\usepackage{tikz}
	\usepackage[all,cmtip]{xy}
	\usepackage{amsthm}
	\usepackage{varioref}
	\usepackage{geometry}
	\geometry{a4paper}
	\usepackage{multicol}
	\usepackage{lmodern}
	\usepackage{hyperref}
	\usepackage{array}
	 \usepackage{fancyhdr}
	 \renewcommand\it\item
\renewcommand{\theenumi}{\alph{enumi})}
	\pagestyle{plain}
	\theoremstyle{plain}
	\fancyfoot[C]{} 
	\fancyhead[L]{Automatismes}
	\fancyhead[R]{Calcul littéral}\geometry{
 a4paper,
 total={210mm,257mm},
 left=10mm,
 top=10mm,
 right = 10mm,
 bottom = 10mm
 }
	
	
	\title{Chapitre 2-  Calcul littéral}
	\date{}
	\begin{document}

\begin{center}{\Large Automatismes Chapitre 2 - Calcul littéral}\\ 
 \end{center}
  \emph{L'objectif des exercices suivants est de travailler vos automatismes. Il est bon de les travailler dès que les méthodes associées auront été vues, et régulièrement ensuite. Pour ce faire, on peut suggérer trois temps : \begin{itemize}
  \item Faire les exercices à l'aide du cours et de la calculatrice (\textbf{N'utilisez pas votre calculatrice pour manipuler les lettres, elle ne le fait pas !}).
  \item Faire les exercices à l'aide du cours sans la calculatrice.
  \item Faire les exercices sans aucune aide.
  \end{itemize}
  Les \textbf{automatismes} ne nécessitent à aucun moment de la réflexion ou de l'initiative, mais uniquement de la méthode. Si vous êtes à l'aise en mathématiques, cela doit donc se faire rapidement et \textbf{sans erreur d'inattention} ; si vous avez des difficultés, ce sont les choses sur lesquelles vous pourrez montrer que vous avez appris sérieusement les méthodes de cours. Il va donc de soi que la notation sera là-dessus binaire : une bonne application de la méthode donne tous les points, la moindre erreur n'en donne aucun. 
  }
  
  \subsection*{Exercice 1 - Évaluer une expression littérale}
  Calculez la valeur numérique des expressions suivantes : 
\begin{enumerate}
\item $7a - 4$ pour $a=3$, $a= -2$, $a=\frac94$.
\it $a(a+1)$ pour $a=1$, $a=-3$, $a=\frac32$.
\it $(2a+1)(a-2)$ pour $a=6$, $a=-1$ et $a=\frac45$.
\it $3a^2-2a+5$ pour $a =  +4$, $a=-2$, et $a=-1$.
\it $3x^2-5x+7$ pour $x =  +3$, $x=+5$, et $x=-3$.
\it $4x^3-12x^2-4x+7$ pour $a =  +5$, $a=-3$, et $a=\frac12$.
\it $\frac{7x^2}2- \frac{8x-6}3 + \frac{3x+7}4$ pour $a =  +3$, $-2$, et $+5$.
\it $\frac{7-3x}{12}+\frac{2(x-2)}3 +\frac54$ pour $a =  +4$, $+2$, et $-3$.
\end{enumerate}
  \subsection*{Exercice 2 - Additionner des expressions homogènes en factorisant}
  
Réduire les sommes suivantes : 
\begin{enumerate}
  \begin{multicols}{2}
\it $3a-a + 4a$
\it $41b+ 3b - 7b$
\it $7ac - 42ac + 12 ac$
\it $30 ab^2 - 12 ab^2 - ab^2 - 72ab^2$
\it $\frac23ax-\frac12ax+\frac34ax-\frac56ax$.
\it $-\frac35a^2bx+\frac14a^2bx-\frac72a^2bx+\frac1{10}a^2bx$.
\it $-\frac47a^2b^3x+\frac52a^2b^3x-\frac54a^2b^3x$.
\it $\frac34a^2b^3x^4y-\frac23a^2b^3x^4y+\frac14a^2b^3x^4y$. 
\it $-\frac32 x + \frac54 x-3x^2+\frac{x}6-\frac52 x^2+5+4x^2$.
\end{multicols}
\end{enumerate}
  \subsection*{Exercice 3 - Multiplier deux expressions littérales simples}
  Effectuer les produits suivants : 
\begin{enumerate}
  \begin{multicols}{2}
\it $ab \times ab^2$
\it $(-abc) \times (2bc)$
\it $(-3ac)(-2a^2)$
\it $(3a^2b^3)\left(\frac23ab^5\right)$.
\it $\left(\frac45a^3b^2c\right)\left(-\frac34abc^4\right)$. 
\it $\left(\frac47a^2xy^3\right)\left(-\frac52a^3y^4\right)$.
\it $\left(-\frac34x^2y\right)\left(+\frac35a^3y^5\right)$.
\it $\left(\frac94a^4x^2y^3\right)\left(-\frac43ax^2\right)$. 
\it $\left(\frac{14}3a^2b^3x\right)\left(-\frac67a^2b^5\right)$. 
\it $\left(-\frac72ax^2y\right)\left(-\frac8{15}b^3xy^2\right)\left(\frac5{21}abx^3\right)$.
\it $\left(\frac5{12}a^4b^2x\right)\left(-\frac27ax^2y^3\right)\left(-\frac{14}5b^2xy^4\right)$. 
\end{multicols}
\end{enumerate}
  
  \subsection*{Exercice 4 - Réduire une somme algébrique}
  
\begin{enumerate}
\it $3a + 4b + 12b -19a$
\it $-8a + b - b^2 + a$
\it $7x^2 + 4x - 3 - 12x + 53 + 12x^2 - 1$
\it $12ac + 41a - 32 ca + 12 c - 4 a$
\it $\frac32 x^2+xy+y^2-2xy+\frac{x^2}3-\frac32x^2$.
\it $4a^2-\frac23a-\frac35a^2+\frac13a-5a-\frac2{15}a^2$.
\it $3x^2+\frac45 - \frac53x - 2x^2 -\frac35x^3 + 4 - 2 x^2 + 7x$. 
\it $4x^2 - \frac72 + \frac35x - \frac52x^2 + \frac43 x^3 - 5 + \frac32 x^3 + 7 - 2x$. 
\it $\frac25 a^2b + 3a^3 -4ab^2 + \frac52 a^2b + \frac72b^3 - b^3 + 2ab^2$. 
\end{enumerate}
  
  
  \subsection*{Exercice 5 - Gestion des signes}
 Simplifier les expressions suivantes : 
  \begin{enumerate}
  \begin{multicols}{2}
  \it $-(x+1)$ 
  \it $-(4-x)$
  \it $-(-x+7)$
  \it $-(-3x-3)$
  \it $x+4 - (x+2)$
  \it $x+4 - x+2$
  \it $x+4 - (x+2) - x + 2$
  \it $x+2-(x+2)$ 
  \it $x+2-x+2$
  \it $x+1 - x - 1 $
  \it $-(x+1) - x - 1$ 
  \it $7x+4- (3x+1) + (x+2)$
  \end{multicols}
  \end{enumerate}
  
  
  \subsection*{Exercice 6 - Distributivité simple}
  
Développer et réduire les expressions suivantes : 
  \begin{enumerate}
  \begin{multicols}{2}
  \it $4(x+1)$ 
  \it $-5(4-x)$
  \it $x(-x+7)$
  \it $2x(-3x-3)$
  \it $3(x+4) - (x+2)$
  \it $(x+4) -5( x+1)$
  \it $5(x+4) - 3(x+2) - 2x + 2$
  \it $-3(x+5)-4(x+2)$ 
  \it $x+4(2-x)+2$
  \it $x+x(1 - x) - 2x(x+1) $
  \it $-(x+1)x - x - 1$ 
  \it $7(x+4)- (3x+1) + x(x+2)$
  \end{multicols}
  \end{enumerate}
  \subsection*{Exercice 7 - Factorisation}
  Factoriser les expressions suivantes : 
  \begin{enumerate}
  \begin{multicols}{2}
  \it $4xy + 12x$
  \it $15ac + 3ab + 6a$
  \it $15ac + 3ab + 6a+ 7a^2$
  \it $15ab + 3ab + 6ab^2+ 7a^2b$
  \it $3x+xy - 5xz$
  \it $10xy + 20yz$
  \it $x^2 + x$
  \it $2(x+1) + 3(x+1) + x(x+1)$
  \it $(x+1)(x+3)+(x+3)(2x-1)$
  \it $(x+1)(x+2) + (x+2)(x+3)$
  \it $(x+5)(2x+5) - (x+5)(4x+2)$
  \it $(x+1)^2 - 3x(x+1)$
  \end{multicols}
  \end{enumerate}
  
  \subsection*{Exercice 8 - Distributivité double}  \begin{enumerate}
  \begin{multicols}{2}
  \it $(x+2)(2x+2)$
  \it $(x-3)(3x-1)$
  \it $(x+5)(-x+7)$
  \it $(4x-1)(-3x-5)$
  \it $(-x-5)(12x+1)$
\it $(2x-3y)(4x-2)$.
\it $(2a+3b)(-4a+6b)$.


\it $(-4x+1)(y-3)$. 
\it $(-2a+3b)(a-b)$. 
\it $(-3y-2+5)(x^2-y)$.
\it $(4a^3+ab)(a^2-b)$.
\it $(5xy+3x)(2x-y)$.
\it $(-3xy-2y)(x+5)$.
\end{multicols}
  \end{enumerate}
  
  \subsection*{Exercice 9 (Pour se préparer à la Seconde) - Distributivité générale }
  \begin{enumerate}
  \begin{multicols}{2}
\it $(-4x+3y+1)(y-3)$. 
\it $(-2a+3b-5)(a-b)$. 
\it $(2x^3-3y-2+5)(x^2-y)$.
\it $(4a^3-5b^4+ab)(a^2-b)$.
\it $(5xy+3x-2y)(2x-y)$.
\it $(-3xy+4x-2y)(x+5)$.
\it $(14a^2b+5a^2-b)(a^2-2b)$.
\it $(7a^3b-4b^2+2a^3)(2a^3+4b^2)$. 
\end{multicols}
  \end{enumerate}
 	\end{document}
